\subsection{From an authored method to a sequence of company statements}
\label{sec:authored-program-composition}

A credit sale creates a claim; its collection later needs that claim to exist.
The complete method text does not by itself determine when a method may be
called or which method version a command selects. The retained registry version
2 still routes its calls through transition version 1, which supports fixed
scalar replacement. Its version number therefore cannot be used to infer
support for claim insertion.

The separate \code{AuthoredMethodRegistry/1} records the transition version
explicitly for each entry, along with the exact authored source, its hash and
the complete method definition and hash. Lowering that source must reproduce
the same method and version. Event-input types, account classes, metadata
bindings and the direct-event-date clock remain explicit. A transition-v2
entry also supplies its memory type definition. The registry validates every
entry against the current book, including entries not selected for the call;
it then dispatches the selected event through the original transition engine.

This state dependence matters when composing commands. A collection method
that reads an identified claim cannot be treated as usable before the sale
creates that claim. \code{AuthoredCompanyProgram/1} therefore supplies a finite
population of explicitly hashed registries and lets each event name the one
it uses. The program first checks its complete command structure and source
identities. At each event, the selected registry is validated against the book
that actually exists at that point. Earlier commands never receive the benefit
of a future state merely to make validation succeed.

Consider the constructed lifecycle used in the preceding claim examples.
The sale increases receivables by 30 and creates the claim. A first receipt of
10 raises cash from 100 to 110 and leaves receivables and the claim at 20.
Due-date assessments change the declared claim observations without creating
cash. The annual report and period carry preserve that remaining claim while
moving the year's profit into retained earnings under the supplied closing
method. A receipt of 20 in the next year brings cash to 130 and clears both
receivables and the claim. Next-period operating cash flow is 20. Two company
charts reproduce this eleven-command sequence, including complete journals,
memory, reports and the period-carry result.

Report and carry commands call the existing native methods with explicit
cutoffs, mappings, equity movements, notes and per-leaf memory policies. The
complete requests and results agree with the retained lifecycle, not only its
cash totals. Method-input schemas remain distinct resources from the schemas
supplied to those report requests. Every successful step also preserves the
typed memory definition. Observing a book at a cutoff does not establish that
all economically relevant events have been supplied.

The comparison has an important representational qualification. The original
lifecycle builder specializes claim identity, customer and due date into its
generated methods, while retaining redundant event fields. The registry
requires exactly the declared inputs, so the derived comparison records every
removed undeclared field and the changed request identity. Complete states,
journals and flows still match the original lifecycle. Exact event-request
comparison is against direct execution of the explicitly derived request.
These methods do not yet supply arbitrary dynamic claim selection.

Failures preserve what happened earlier. If the first sale succeeds and the
next collection amount is unknown, the result retains the sale, its claim and
the unexecuted command suffix, with the collection unresolved. Negative amounts,
false applicability, repeated identities, early reporting cutoffs, incomplete
carry policies and invalid unselected entries retain their distinct refusals.
Seven such cases preserve the exact successful prefix. Forged complete results
fail recomputation. A separate authored edit halves the collection algorithm's
payment: after the sale and that receipt, cash is 105 and receivables and the
claim are 25. This intervention tests execution of the supplied algorithm.

The checked registry and program remain separate reference successors. They
preserve the original engines and their provisional status. The program's
64-command and eight-mebibyte trace limits are reference resource bounds,
not production capacity claims. Full textual authoring of registry resources,
general method specialization, broader financial methods and actual-company
source admission remain open. The comparisons are retained under
\path{data/accounting-master/authored-method-registry-20261026/} and
\path{data/accounting-master/authored-program-composition-20261026/}.

\subsection{A typed claim can still disagree with the book}
\label{sec:authored-state-invariants}

The preceding successful histories do not establish that every authored method
preserves the original lifecycle's economic controls. Suppose an otherwise
unchanged sale posts receivables of 30 but creates a claim whose remaining
amount is 31. Both values have the right monetary type, the journal balances,
and the memory record satisfies its schema. The original lifecycle nevertheless
rejects that state because the claims do not sum to receivables. Typed-book
validation alone does not express this cross-record relation.

A focused audit changes one claim field at a time in the authored sale and
then requests a report from the genuine period opening. Across two company
charts, sixteen cases expose eight outer-control gaps: claim totals, the sale
event's membership in the book, issue/due chronology, future issue dates,
negative assessment amounts, nonzero unassessed placeholders, future assessment
and an incorrect overdue snapshot. An empty customer string is already refused
by the memory schema. That covered structural rule must not be counted as a
missing economic invariant.

\code{BookStateInvariant/1} supplies the missing connection to the existing
collection evaluators. Its explicit bindings select the book's date and
receivable account, the identified rows of its claim map and its seen-event
identities. Each scalar's type and context must agree with the declared input,
including for an empty map. Values come from that exact supplied book; a caller
cannot substitute a Boolean assertion that its controls passed.

The unchanged collection evaluator checks unique claim identities, references
from each claim to its sale event and guarded conditions on every row. The
unchanged collection-total evaluator sums remaining amounts and compares that
total with receivables. An assessed claim must have an assessment date between
issue and the current date. Its overdue snapshot equals the assessed amount
when assessment occurs strictly after the due date, and zero otherwise.
An unassessed claim retains zero assessment and overdue placeholders. Remaining,
assessed and overdue amounts are nonnegative. These are declared controls of
this constructed lifecycle, not a new impairment model.

The separate \code{GuardedAuthoredCompanyProgram/1} selects the invariant by
its exact hash. It checks the opening, each current book and each proposed
after-state. Only a proposed state that satisfies the selected controls enters
the successful prefix. In the 31-versus-30 example, the proposed sale is
refused, the opening book remains current, and the sale and report remain in
the unexecuted suffix. A failed post-state check cannot quietly retain the
invalid claim or its otherwise balanced journal.

Both eleven-command histories retain their complete original results, and
all 24 opening and after-state books pass the original lifecycle control.
All sixteen changed-method counterexamples are refused. Wrong pointers,
contexts and field types also refuse. The original evaluators retain unresolved
results for partial or unknown collection inputs, while an unrelated unknown
does not erase a known contradiction. Unsupported capabilities and explicitly
unsupported consequences remain distinct from invalid states. Neither a
successful invariant nor a checked schema establishes source or accounting
admission.

Completeness here means all members of the supplied internal map. It does not
mean that an actual issuer disclosed every claim. Event-specific duties also
remain: selecting the right customer, preventing an inappropriate new event
and respecting sealed assessment windows require their own prerequisites.
The state-control repair therefore preserves the finite comparisons without
claiming whole-lifecycle substitutability. Counterexamples and repair evidence
are retained under
\path{data/accounting-master/authored-program-invariant-audit-20261026/} and
\path{data/accounting-master/authored-state-invariants-20261026/}.

\subsection{Keep the event's intent when a method is specialized}
\label{sec:authored-event-contract}

A consistent closing book still cannot show whether a collection was permitted
at the time it occurred. In the constructed lifecycle, an assessment on
31 January closes a window and advances its cursor to 1 February. A collection
dated 31 January can satisfy ordinary nondecreasing book dates and all claim
state invariants, yet violate that sealed-window policy. The retained lifecycle
refuses it. A separately selected event contract is needed to preserve that
meaning when the method is authored elsewhere.

Customer and interval assertions expose another boundary. The original full
event states which claim and customer a receipt concerns, or which start and
end dates an assessment covers. A specialized method may embed those values,
so they no longer appear among its declared calculation inputs. Removing them
from the projected event is acceptable only if the original intent remains
available wherever eligibility and specialization must be checked. A method
that happens to produce a consistent state cannot supply the missing assertion
about which customer the caller meant.

\code{ContractedAuthoredCompanyProgram/1} therefore retains every complete
original event beside the projected event. Its caller explicitly selects the
retained receivable lifecycle contract, identified by the native implementation
hash, roles, unit-cost pointer and policy. Before an event, that named owner
checks the full input population, unknowns and applicability on the actual
current book, then derives the intended method under its existing claim,
customer and interval rules. The complete authored materialized request must
equal that intended request after the declared input projection and permitted
scalar-collection transition conversion. Both requests and the removed input
values remain inspectable.

This contract is deliberately specific. The half-payment method previously
authored is a valid alternative algorithm within its stated scope. It does not
satisfy a caller's selection of the original full-payment lifecycle, so that
selection refuses it. An explicit policy choice resolves the difference; a
balanced journal or a valid type cannot choose between the two algorithms.
The native specialization owner is named, not hidden behind an arbitrary
callback or presented as a general declarative financial compiler.

Both eleven-command examples preserve complete states, statements, period
carry and archives, and every full-intent native event request matches the
retained lifecycle. Eighteen refusal comparisons across the two charts cover
customer mismatch, missing or repeated claims, unknown inputs, false
applicability, inconsistent assessment intervals, an early due date and the
sealed-window event. Earlier successful commands remain unchanged when a later
event refuses. Changed projected values and a method inconsistent with the
selected intent are also refused; forged intent evidence fails recomputation.

These comparisons extend the checked refusal scope without proving every
lifecycle rule or supplying actual company facts. General portable authoring
of event contracts, wider financial operations and source/accounting admission
remain unresolved.

The source boundary and counterexamples are retained under
\path{data/accounting-master/authored-event-prerequisite-audit-20261026/}.
The selected-contract evidence is retained under
\path{data/accounting-master/authored-event-contract-20261026/}.

\par

\subsection{Constructing a method from typed intent and ordered claim rows}
\label{sec:method-specialization}

A sale method must insert the claim the caller actually names. The same
financial calculation can apply to a new customer, due date or invoice key,
yet the inserted record and the later collection pointer must change. Writing
one finished method for each saved event would merely move the native
generator's results into a catalogue. Explicit parameters instead retain the
complete method algorithm and identify exactly which values come from intent
or current state.

\code{ParameterizedCompanyMethod/1} supplies that bounded construction step.
Its template consists of literal, object, array and named-parameter nodes.
Each parameter has a declared scalar type and source: a known full-intent input,
event identifier or date, selected current-book metadata, or a pointer assembled
from literal segments and earlier text parameters. A claim key containing a
slash or tilde is escaped as a pointer segment; it is never reparsed as source.
Unknown required input remains unresolved. Duplicate, unused, forward or
unbound parameters are refused.

For a sale, parameters supply the claim key, customer, due date, event identifier
and dates in the inserted row. For a collection, they supply the selected
claim key and its remaining-balance pointer in both the precondition and the
monetary input binding. The subtraction, postings and memory write remain
visible in the full method. The input schema is compiled from its explicit
typed definition, while the state schema comes from the actual book. Complete
template, state and event identities bind the constructed request. The existing
typed frontend and method-source renderer then check the whole method.

Four templates across the two retained charts reproduce six prior sale and
collection requests and twelve new events with changed keys, customers, dates
and amounts. Comparisons cover complete requests and transition results, and
the emitted method source compiles back exactly. A deliberate edit halves the
collection algorithm's cash, receivable, flow and remaining-balance effects.
It produces exactly half the cash movement and preserves the claim balance
relation, while differing from the explicitly selected full-payment contract.
The construction mechanism therefore exposes a consequential algorithm rather
than hiding a table of fixed financial outcomes.

Assessment adds a different need: the number of writes depends on the current
claim population. \code{OrderedMethodExpansion/1} declares bounded row
expressions and an ordered list of claim keys. In the retained state-only
assessment, each row sets the assessment flag and date, copies remaining
amount into the snapshot, and records overdue amount only when the due date
is strictly earlier than the assessment date. A final write advances the
cursor. Typed expressions can read a row field, key, event metadata or known
event input, compare dates and choose between equally typed values. Both
conditional branches and every write destination are checked even when the
claim map is empty.

Order is explicit because it affects construction identity. Canonical JSON
hashing ignores the order of object fields, but the native assessment's ordered
write list follows its claim-map iteration. The new call therefore binds a
separate key-order digest and requires the list to cover every current claim
exactly once. Forward and reverse orders produce their corresponding native
requests. Their disjoint writes yield the same financial state, but that
agreement does not make the two ordered request records identical.

Eight retained and twelve new assessments match full requests and transition
results, including empty maps, multiple claims, escaped keys and due-date
boundaries. An authored inclusive comparison makes a claim overdue on its due
date. The typed transition executes that different rule, while the selected
strict lifecycle invariant rejects its snapshot. This distinction matters:
types can establish that a date comparison is meaningful, but they cannot
choose the caller's required convention.

Both constructors work with the native specializer disabled in their checks.
Their construction interface leaves full-intent eligibility separate; the next
section supplies an explicit ordered alternative for the retained contracts.
Fixed chart context, role mapping and unit-cost choices remain declared scope.
The earlier complete program composition and invariant source continue to
provide their own boundaries.

\begingroup\raggedright
Complete constructions, controls and variants are retained under
\path{data/accounting-master/parameterized-method-construction-20261026/}
and \path{data/accounting-master/ordered-assessment-expansion-20261026/}.
\par
\endgroup


\subsection{Preserving the order of eligibility decisions}
\label{sec:ordered-event-eligibility}

A customer can be unknown before a collection is authorised. Suppose the full
intent supplies an unknown customer, false approval and a claim key absent
from the book. The retained lifecycle first reports unresolved input. Once the
customer is supplied, the same false approval produces an unsupported event;
only after approval can the missing claim produce an invalid event. Combining
all three predicates into one conjunction changes this behaviour: an invalid
consequence can outrank the earlier unresolved or unsupported result. The
individual predicates remain meaningful, but their conjunction computes a
different decision.

\code{OrderedEventEligibility/1} makes this order explicit. A template declares
its company context, operation, complete input types and ordered stages. Each
stage supplies one existing guarded requirements manifest, exact bindings for
its declared inputs, and the reason to report if it does not pass. The evaluator
validates the typed book, event identity, chronology, full input population,
types, values and provenance before interpreting the stages. Any unknown full
input then returns unresolved, preserving the input names. Approval is the
first semantic stage in each retained template. A stage with no active
requirement is refused rather than treated as permission to continue.

For a sale or collection, the next stage checks that assessment has not sealed
the event date. A sale then checks the due date and absence of the claim key.
A collection checks presence of the key before comparing its recorded customer
with the supplied customer. An assessment instead checks its starting cursor,
nonempty interval and last included date, then its period boundary. Execution
stops at the first stage that does not pass, retaining its status, reason and
stage identifier. The existing guarded evaluator still determines each typed
predicate; the ordered layer introduces no competing Boolean language.

The bindings supply facts rather than conclusions. They can read a full-intent
input, event identifier or date, book date or period end, a declared memory
leaf, map membership, or a selected row's scalar field. Membership is computed
against the actual typed map using the supplied text key. A caller cannot
replace that computation with an asserted Boolean. A row-field binding names
an earlier stage whose unconditional predicate establishes membership in the
same map under the same key. Compilation checks that relationship and the
declared field type without requiring a row to exist. It also checks every
later stage and guarded branch before execution, including when the map is
empty or an earlier approval will fail.

Calendar adjacency is a narrow additional binding. After strict date
validation, it compares the two date ordinals and returns true when the right
date is one calendar day after the left. An assessment therefore relates its
event date to its explicit end-exclusive date without constructing a successor
that could overflow at the maximum representable date. Leap days and year
changes follow the same relation. This is a civil-calendar condition, not a
business-day adjustment or a fiscal-calendar policy. The unchanged chronology
and interval checks remain separate conditions.

The call binds the template, book, event, schema registry and memory definition.
It also binds the order of full-intent inputs. Canonical JSON identity ignores
map insertion order, whereas the retained unresolved result lists unknown
inputs in that order. Reversing two unknown inputs therefore leaves the
canonical event hash unchanged but changes the order digest and the returned
list. Reusing the earlier call is refused. This is the same distinction between
value identity and explicitly meaningful order that arose for assessment writes.

Six templates across the two retained charts admit all 38 previously checked
events. Combined with the parameter and ordered-row constructors, they reproduce
the complete saved requests and transitions while native lifecycle eligibility
and construction functions are disabled. The original 18 refusals and 16
additional precedence cases preserve their statuses and reasons. Seventy
controls include selected-row binding, skipped-stage compilation, changed rule
truth, stale identities and calendar boundaries. A changed customer predicate
can admit a wrong customer; that consequential edit demonstrates author control
and deliberately differs from the selected contract.

These comparisons establish a finite declarative alternative for eligibility
and method construction. Complete-program composition adds a different question:
do reports, carry and invariant failures preserve the same current state and
successful prefix? The following extension checks that execution boundary
while retaining opening admission as a separate prerequisite. General company
rules and independent implementation acceptance remain open.

\begingroup\raggedright
The ordered rules, full comparisons, identity repair and control results are
retained under
\path{data/accounting-master/ordered-event-eligibility-20261026/}.
\par
\endgroup


\subsection{Executing declared event rules through a company program}
\label{sec:declarative-event-program}

A collection that works in isolation can still fail in a company program.
An earlier sale changes the claim population and its state schema; a report
advances the observation date; a year-end carry establishes the next opening
book. The method for a later event must be constructed from that current state.
Replaying a previously finished method cannot establish that these dependencies
are handled correctly.

\code{DeclarativeEventCompanyProgram/1} composes the explicit mechanisms.
The program supplies its opening book, typed memory definition, report mapping,
authored state invariant and one eligibility/construction pair per operation.
Event commands contain the full intent. An assessment also supplies an ordered
list covering its current claim population. The program contains neither native
event callbacks nor precomputed event results. Separate hashes bind the selected
rules and invariant; program identity also retains the order of event inputs.

Before each command, the runner checks identity uniqueness, the typed book and
the selected state invariant. For an event, ordered eligibility runs first.
Only an eligible event reaches the parameter or ordered-row constructor, which
uses the current book and full intent to produce the complete transition
request. The unchanged typed transition then computes postings, flows and
memory effects. A report uses the current period's opening and the book observed
at its stated cutoff. A carry uses the existing closing and carry evaluator,
with the explicitly selected policy preserving every memory leaf.

The book advances only after the resulting state also passes its type and
invariant checks. Consider an authored sale whose postings increase receivables
but whose inserted claim incorrectly records zero remaining amount. The typed
transition can propose that state because each value and destination has a
valid type. The separately selected claim-total invariant rejects it: claim
amounts no longer equal the receivables balance. The runner retains the prior
current and opening books, records the refused command, and leaves the complete
remaining suffix unexecuted. It adds no compensating journal or financing event.
This is in-memory refusal behaviour, not a claim of durable database rollback.

The two retained eleven-command programs preserve all complete native command
requests and results, including reports and carry archives, along with their
before/after state identities and final books. Two further eleven-command
programs start from the same declared openings but create three new claims each,
using changed keys, customers and amounts. Their report, year-end carry and
following-year assessment also agree. All native lifecycle functions are
disabled during candidate execution. The saved result is checked again after
serialization; no native after-state is supplied to the candidate.

Twenty-six adverse programs include the original eighteen refusal cases,
combined unknown/approval/missing-claim conditions, and negative or excessive
collections that fail financial preconditions after eligibility has passed.
The latter retain the complete refused request and transition. Additional
controls check stale selections, changed input order, duplicate identities,
incomplete assessment order, a changed carry policy and the invariant-invalid
sale. Each stopped run preserves its successful prefix, current opening,
reports and unexecuted suffix.

The comparison concerns complete substantive command payloads. The candidate
adds explicit eligibility, construction and invariant records; the native
result carries different implementation and observation records. Those
differences are retained and identified, rather than removed and described as
whole-result equality. Native observations are separately recomputed from the
candidate states. Shared transition, report and carry evaluators remain the
owners of their respective financial calculations.

Opening admission remains a separate native prerequisite for this composer
comparison. The original validation checks the retained opening/program before
native functions are disabled. The next extension declares those opening
conditions and binds them to the exact program selected for execution. The
earlier company-source frontend retains its own native-contract target.
General opening populations, a complete financial language, accounting validity
and independent implementation acceptance remain separate questions.

\begingroup\raggedright
The four complete programs, adverse cases and full comparisons are retained
under \path{data/accounting-master/declarative-event-program-20261026/}.
\par
\endgroup


\subsection{Admitting the opening that the selected program will use}
\label{sec:declared-opening-binding}

A balanced, typed opening book can still violate the selected company's
operating conventions. Its claim map may be empty while receivables are
nonzero; its unit-cost leaf may have the wrong dimensions; its account roles
may overlap; or its statement mapping may omit an account. The opening check
must establish these additional conditions before events run. Yet it must not
silently promise that every future command is valid: an unknown future customer
or an unsupported event is still decided at its execution boundary.

\code{DeclaredOpeningAdmission/1} records the retained empty-opening contract
as ordered declarations. Its template specifies the input envelope, policy
fields and values, exact claim-row and clock types, empty claim population,
cursor equal to opening date, complete and distinct account roles, required
account classes and cash identity. It then checks the unit-cost pointer,
dimensions and strictly positive value, the opening receivable total, the
opening statement report and the command-list limit. Each condition retains
its native position and failure reason. All template stages compile before
their evaluation; a malformed later declaration cannot disappear because an
earlier condition would refuse the opening.

Structural conditions compare declared record fields, schema nodes, role maps
and literal policy values. Scalar date and monetary predicates use the existing
guarded requirements evaluator, retaining their dimensional types until that
evaluator lowers them. This separation preserves a subtle input boundary:
a malformed carry-policy value fails the declared policy equality, rather than
acquiring an earlier and different scalar-type error. Types remain essential
where the selected contract requires typed numerical or date comparisons.

The empty-claim premise makes the opening claim-total condition particularly
simple: the receivables balance must be zero. This is not a replacement for
the row invariants of an operating book. Once events create claims, the complete
authored state invariant again checks their totals, dates, references and
assessment snapshots. The opening report retains its own evaluator, which
checks the period start, opening flows and temporary balances, statement-map
coverage and statement arithmetic. No statement rule is duplicated inside the
opening template.

Two templates reproduce exact acceptance or refusal on 82 complete inputs,
including the 58 preceding audit cases. Ten are admitted. Among them are
inputs with a malformed future command, unknown operation or unknown customer
that the original opening validator also accepts. The corresponding execution
checks remain separate. A deliberate change from strictly positive to
nonnegative unit cost admits a zero-cost opening; this proves that the declared
predicate controls the decision and explicitly differs from the selected native
contract. Dependency identities, input order and stale selections are checked
alongside the substantive conditions.

Passing opening admission and passing execution independently would still
leave a serious gap. One could admit a book under one account-role map but
execute methods selected for another, or admit a different command list from
the one that actually runs. \code{AdmittedDeclarativeProgram/1} therefore binds
both evaluations to one explicit selection. That selection names the chart
context, role map, unit-cost pointer, policy, opening template, complete rule set
and state invariant. The call binds all supplied definitions and facts rather
than accepting a caller's success flag.

After admission, the adapter compares the complete shared opening, schemas,
statement mapping and memory definition. It checks the selected rules and
invariant, then compares every projected event, report and carry command.
Projection removes only the assessment-key order and memory carry-policy fields
that the native-shaped admission input lacks; the composer checks those fields
when they are used. Event-input order is retained explicitly because it affects
the ordered list of unknown inputs. A mismatch refuses execution.

With all native lifecycle functions and the operating-method builder disabled,
the combined path preserves the complete prior composer results for four
successful programs and 26 adverse programs. The successful programs contain
44 commands whose complete substantive payloads already agree with the native
baseline. Twenty-two binding controls reject changed facts, roles, cost
selection, policy, rules, invariant, commands, order and stale identities.
An opening refusal executes no command. Existing typed-memory, transition,
requirements, report and carry evaluators remain the owners of their respective
calculations.

This removes the native opening evaluator prerequisite for the selected
combined path. The selection is still a finite reviewed chart contract, not
a proof that arbitrary edited rules have equivalent meanings. The portable
inputs use exact JSON values and rational strings. Earlier company-source
frontends retain their older target. Section~\ref{sec:selected-company-source}
introduces the successor that expresses these algorithms directly. General nonempty opening populations, wider company rules,
independent implementation and accounting acceptance remain open.

\begingroup\raggedright
The opening comparisons and bound programs are retained under
\path{data/accounting-master/declared-opening-admission-20261026/}
and \path{data/accounting-master/admitted-program-binding-20261026/}.
\par
\endgroup


\subsection{Compose methods, state controls, and event intent in one authored program}
\label{sec:common-program-authoring}

The preceding examples established four separate boundaries: a method can be
written in complete text, a registry can retain its source and identity, a
state invariant can reject a bad book, and a selected event contract can retain
the caller's full intent. A usable program must express those choices together.
Otherwise a typed fixture may pass while the proposed language has no way to
write the same program or to show which semantic owner supplied a missing field.

\code{AuthoredCompanyText/1} closes this composition boundary with a bounded
source format. The program imports the opening book, report schemas, statement
mapping and memory definition by exact SHA-256 identities. It selects one
\code{BookStateInvariant/1} resource explicitly. These imports are facts and
checked rule resources; they do not silently generate methods or claim that a
small synthetic state is a complete disclosure history.

Each registry entry contains the complete
\code{CompanyMethodAuthoring/1} source block. The entry records its lowered
method digest, source digest, transition profile, input declarations, account
classes, clock and separate entry-schema resources. A registry version therefore
cannot hide a transition change. The method header supplies the transition
profile, while the program author chooses the registry and method for each
event. Entry schemas remain distinct from the shared report schemas because
combining them would change the identity of a report request.

For the selected receivable contract, the same program also declares the native
implementation identity, role-to-account map, unit-cost pointer and policy.
Every event carries both its projected method inputs and its complete original
intent. The latter retains the customer, claim and interval assertions that a
specialized method may remove from its calculation inputs. The runtime then
checks the selected native lifecycle on the current book and requires the
authored projected request to equal the declared projection of that full
request. An edited half-collection method therefore remains executable under
the guarded target, where it changes the resulting claim and cash, but it is
refused when the caller explicitly selected the full-payment native contract.
This is a contract comparison, not an inference that either method is the
correct accounting treatment.

Reports, notes, carry boundaries and memory policies reuse the existing command
forms. A reset policy can obtain its type only from an explicitly present
opening memory leaf. Unknown event inputs retain their runtime unresolved
status. Duplicate declarations, unregistered aliases, changed resources and
incomplete source blocks are errors with a stage and UTF-8 byte location.
Method source is hashed from its original bytes; the
verifier checks outer, method and declaration spans after Unicode and CRLF
roundtrips.

Two complete eleven-command programs reconstruct their saved typed programs
and runtime results exactly. The run includes nineteen structural refusals and
six matched native refusals across the two charts, direct method editing,
report-note preservation and failed-prefix checks. These finite comparisons
show that complete methods, explicit registries, selected state rules and full
event intent can share one source boundary. They do not provide general
invariant authoring, whole-company source completeness, accounting admission,
production readiness or human acceptance. The following successor supplies
explicit source for the current finite invariant family.

\begingroup\raggedright
The field inventory and grammar design are retained under
\path{data/accounting-master/common-program-authoring-audit-20261026/}.
The implementation and checks are under
\path{data/accounting-master/common-program-authoring-20261026/}.

\par
\endgroup

\input{sections/invariant-source-authoring}

\input{sections/selected-company-source}

\input{sections/complete-claims-opening}

\input{sections/nested-term-invariants}

\input{sections/selected-term-execution}

\input{sections/selected-financing-history}



